% !TeX root = ../main.tex
% High-stress and time-dependent ground behaviour. Moved to Part III:
% squeezing, swelling, spalling and bursting are behaviour types, and
% belong with the ground behaviour taxonomy rather than with design.
\chapter{Squeezing, Swelling and High-Stress Ground}
\label{ch:high-stress}

Chapter~\ref{ch:ground-behaviour} classified ground by the mechanism that governs its response to excavation and treated the cases in which that response is immediate. This chapter takes three behaviours that need separate treatment: ground that closes on the excavation over months, ground that expands on contact with water, and strong rock that fails violently when its stored energy is released.

Squeezing and bursting share a design consequence: a stronger structure is rarely the answer. A tunnel survives squeezing by yielding in a controlled way and survives a burst by absorbing energy. Stiffness in the support, the usual response to a structural problem, makes both worse. Swelling differs, because the choice there lies between resisting the pressure and allowing the strain.

\section{Squeezing and swelling ground}
\label{sec:squeezing-swelling}

\subsection{Squeezing mechanism and prediction}
\label{sec:squeezing-prediction}

Squeezing is time-dependent large deformation occurring where the rock mass
strength is low relative to the in-situ stress. The ground does not collapse; it
closes on the excavation progressively, and the closure
continues after the face has passed.

Prediction is empirical and reasonably reliable. The controlling parameter is
the ratio of rock mass strength to the mean in-situ stress at tunnel depth
(Chapter~\ref{ch:in-situ-stress}). The strain chart of
\citet{hoekmarinos2000squeezing} relates that ratio to the percentage closure to be expected,
in five classes. These run from below about one per cent, requiring no special
measures, to more than ten per cent, where extreme squeezing problems arise
and conventional support cannot succeed. The relation is set out in
Section~\ref{sec:screening-squeezing} of Chapter~\ref{ch:empirical-design},
together with the definition of rock mass strength on which the chart is
calibrated, and the class boundaries are in
Table~\ref{tab:squeezing-classification} of that section.

The prediction is available at design stage from
quantities already in hand (the intact strength, $m_i$ and GSI), and
the class in which the predicted strain falls indicates the strategy
required (Section~\ref{sec:squeezing-strategy}).

Time dependence must be separated into two mechanisms that look alike in the
data. Some convergence is face-advance deformation, the
three-dimensional stress redistribution captured by the longitudinal
deformation profile (Section~\ref{sec:ldp-appendix}), and it would occur in
a rate-independent material. The
remainder is creep, continuing at constant face position. Only the
second requires a viscoplastic model, a constitutive law in which strain
continues to accumulate with time under constant stress once a yield condition
is exceeded (Section~\ref{sec:scope-limits-numerical}). Attributing all observed
time-dependent convergence to creep overstates the problem.

\subsection{Support strategy and deformation allowance}
\label{sec:squeezing-strategy}

Three strategies exist. \emph{Resist} installs support stiff and strong enough
to prevent significant convergence. It is viable only for mild squeezing, that
is the minor squeezing class of
Table~\ref{tab:squeezing-classification} in Chapter~\ref{ch:empirical-design},
and it rapidly becomes uneconomic because the required pressure rises steeply as
convergence is suppressed. \emph{Yield} allows controlled convergence through
sliding sets, deformation slots and yielding bolts, so that equilibrium is
reached at a support pressure the system can carry
(Chapter~\ref{ch:ccm}; the construction that compares a yielding system against a stiffer one is set out in Section~\ref{sec:seq-support-interaction} of Appendix~\ref{app:calculations}).
It is the normal answer in the severe class, with the over-excavation
allowance set out below, and in the very severe and extreme classes it is
combined with reprofiling planned from the outset. \emph{Avoid} changes the alignment or depth to bypass
the squeezing ground altogether, and where the squeezing zone is localised
this is often the cheapest answer, above all in those two classes.

Yielding requires the profile to be over-excavated. The finished profile must
survive the convergence, so the excavated profile is enlarged by the
anticipated closure plus a margin. Getting that allowance wrong in
either direction is costly, since too little means reprofiling and too much
means excavating and lining a larger tunnel than required. In extreme cases
reprofiling is planned and repeated as a normal part of the cycle.

The allowance is a radial dimension: the predicted closure strain multiplied
by the radius, with a margin for the scatter of the prediction. For the poor
rock of the deep tunnel of Chapter~\ref{ch:ccm}, with a radius of 4.75~m,
the mean prediction of 3.7 per cent, from the strain relation of
Section~\ref{sec:screening-squeezing}, gives 176~mm, and the 84th-percentile
strain of 6.5 per cent obtained in
Section~\ref{sec:reliability-support-uncertainty} gives about 310~mm
(Figure~\ref{fig:squeezing-allowance}). The
difference between the two is what the margin has to cover,
and it spans a class boundary: the mean lies in the severe class and the
84th percentile in the very severe one.

\begin{figure}[htbp]
\centering
\datafigure[\textwidth]{ch14/squeezing-allowance}
\caption{Closure strain against GSI for the poor rock of Chapter~\ref{ch:ccm} ($\sigma_{ci} = 72$~MPa, $m_i = 7$, $p_0 = 16.9$~MPa), with the radial allowance $\varepsilon a$ for $a = 4.75$~m on the right axis. The grey band is GSI $30 \pm 5$; the mean and 84th-percentile strains fall in the severe and very severe classes.}
\label{fig:squeezing-allowance}
\end{figure}

The contractor, however, sets out and surveys a profile, not a radius. The
radial allowance therefore has to be converted into an enlarged design
profile, with the survey datum against which overbreak and underbreak are
measured agreed, before excavation starts; on a non-circular profile the
predicted closure, and with it the allowance, differs between crown, wall
and invert.

\subsection{Swelling ground}
\label{sec:swelling-ground}

Swelling is a different mechanism producing similar symptoms. Smectite-rich
mudrocks swell on contact with water; anhydrite converts to gypsum with a
large volume increase. Their identification at investigation stage is
covered in Chapter~\ref{ch:engineering-geology}. Both generate pressure that continues for years and is
concentrated at the invert, where water collects. The distinction from
squeezing matters because the remedies differ: swelling responds to keeping
water away from the rock, and where it cannot be prevented the choice is
between resisting the pressure with a heavy invert or accommodating it with a
compressible layer. Swelling pressures are measured in oedometer swelling
tests on core. Indicative magnitudes are 0.1 to 2~MPa for smectite-bearing
mudrocks and several megapascals for anhydrite-bearing rock, and it is
towards the upper end of that range that a compressible layer becomes the
practical choice.

\section{Brittle failure and rockburst}
\label{sec:brittle-rockburst}

\subsection{Brittle failure at high stress}
\label{sec:brittle-failure}

Strong brittle rock at high stress fails differently from weak rock. Rather
than yielding progressively, it spalls: thin slabs detach from the boundary,
progressively deepening a characteristic notch where the tangential stress is
highest. The process is driven by extensile cracking parallel to the boundary
rather than by shear.

Conventional criteria predict this poorly.
Near a free surface, confinement is close to zero, and in that condition
cohesion is lost at small strain while friction requires larger strain to
mobilise, so the two are not available simultaneously. Criteria that add
them overestimate near-boundary strength. Formulations in which cohesion
degrades and friction mobilises with accumulating strain reproduce that
sequence, and so predict the lower strength near the boundary.

The practical prediction is empirical. \citet{martin1999spalling} found
that the depth of the spalling notch grows linearly with the ratio of maximum
boundary stress to unconfined compressive strength, with spalling initiating
at a ratio of roughly 0.4, far below the strength the rock exhibits in a
laboratory test (Figure~\ref{fig:spalling-notch}). The relation is given in
Section~\ref{sec:screening-brittle} of Chapter~\ref{ch:empirical-design},
where Table~\ref{tab:spalling-example} works the case drawn in the figure.
It is the basis for setting reinforcement length
in spalling ground, within the limits stated with it in that section: massive to moderately jointed strong
rock in which failure is stress-driven, and not weak or closely jointed rock
masses. Along the tunnel the notches are a pair of V-shaped grooves parallel
to the axis, on the sides of the opening normal to $\sigma_1$
(Figure~\ref{fig:r3d-spalling-notch}).

\begin{figure}[htbp]
\centering
\datafigure[\textwidth]{ch14/spalling-notch}
\caption{Spalling notch for the case of Table~\ref{tab:spalling-example} ($\sigma_1 = 60$~MPa, $\sigma_3 = 11$~MPa, UCS 224~MPa). Left: the notch, drawn to the depth $r = 1.44a$ of the relation over the arc where the tangential stress exceeds 0.4\,UCS (notch shape schematic). Right: tangential stress on the boundary, above 0.4\,UCS within $39.5^\circ$ either side of the normal to $\sigma_1$.}
\label{fig:spalling-notch}
\end{figure}

\begin{figure}[htbp]
\centering
\renderfigure{spalling_notch}{%
  \rlab[south]{0.33,0.70}{0.22,0.90}{Spalling notch}
  \rlab[north]{0.40,0.33}{0.33,0.08}{Slabs}
}
\caption{The notches of Figure~\ref{fig:spalling-notch} along a tunnel three diameters long, with a wedge cut away about the right springline. With $\sigma_1$ horizontal they form at crown and invert, to $r = 1.44a$ over $\pm 39.5^\circ$. The carmine skin marks the spalled flanks, and slabs lie in the invert notch (radius and slabs illustrative; notch shape schematic).}
\label{fig:r3d-spalling-notch}
\end{figure}

\subsection{Energy and excavation stability}
\label{sec:energy-stability}

Excavation releases stored strain energy (Chapter~\ref{ch:rock-mechanics}). Part is absorbed in deformation of
the surrounding rock, part does work on the support, and the remainder
(excess energy) is available to drive failure. At moderate depth it
dissipates without violence; at high stress in brittle rock it may not.

The energy release rate, the stored strain energy released per unit area of
excavation created and usually quoted in megajoules per square metre, serves
as a design index for deep excavation, and its
value lies in the comparison of layouts and sequences rather than in an
absolute threshold. A sequence that releases energy gradually is preferable to
one that releases the same total abruptly.

Whether failure is stable or violent depends on relative stiffness. If the
surrounding rock mass is stiffer than the post-peak stiffness of the failing
volume, load sheds gradually and failure is controlled. If it is softer, the
surrounding rock unloads faster than the failing volume can shed strength, and
the stored energy drives the failure dynamically (Figure~\ref{fig:burst-stiffness}). This is the same criterion
that distinguishes a yielding pillar from a bursting one
(Chapter~\ref{ch:openings}), and it explains why a stiff loading system is
protective while a soft one is not.

\begin{figure}[htbp]
\centering
\datafigure[0.62\textwidth]{ch14/burst-stiffness}
\caption{Stable and violent failure (schematic). The black curve is the response of the failing volume; the lines from the peak are the unloading of the surrounding rock. A surround stiffer than the post-peak slope sheds load gradually; a softer one releases the shaded excess energy into the failure.}
\label{fig:burst-stiffness}
\end{figure}

\subsection{Rockburst and dynamic support}
\label{sec:rockburst-support}

Three burst types are distinguished by source. A \emph{strainburst} is a
violent failure of the rock at the excavation boundary, driven by locally
stored energy. A \emph{pillar burst} is failure of a highly stressed pillar. A
\emph{fault-slip burst} originates on a structure remote from the excavation
and reaches it as a seismic wave, so the source may lie well outside the
zone of influence of the opening, and the damage appears wherever the wave
meets ground that is unsupported or poorly supported.

Dynamic support is designed in terms of energy. The demand is the kinetic energy of
ejected rock, which per unit area of excavation surface is
\[
E = \tfrac{1}{2}\,\rho\, t\, v^{2},
\]
for a slab of thickness $t$ and density $\rho$ ejected at velocity $v$; where
the slab falls from the back, the work done by gravity over the support's
displacement adds to it. A slab 0.3\,m thick of density
2700\,kg/m\textsuperscript{3} ejected at 3\,m/s carries
3.6\,kJ/m\textsuperscript{2} (Figure~\ref{fig:burst-energy}). The ejection velocity is assumed rather than
measured, commonly at a few metres per second, and because it enters as its
square it is the dominant uncertainty in the method. The capacity is the
energy the support system absorbs
before failing, which requires yielding elements: a stiff, high-capacity bolt
absorbs almost nothing because it accommodates no displacement. Representative
absorption capacities, rated in kilojoules rather than kilonewtons, are in
Table~\ref{tab:support-props} of Appendix~\ref{app:constants}.

System compatibility governs the outcome, and it is the design point most
often missed. A bolt capable of absorbing energy over 200\,mm of displacement
achieves nothing behind a mesh that tears at 50\,mm, or a sprayed concrete (shotcrete) shell that
cracks and sheds at 20\,mm. Every element must be able to deform through the
same displacement, and in burst-prone ground the retention system rather than
the reinforcement is usually the weak link.

\begin{figure}[htbp]
\centering
\datafigure[\textwidth]{ch14/burst-energy}
\caption{Left: kinetic energy of an ejected slab, $E = \tfrac{1}{2}\rho t v^2$ with $\rho = 2700$~kg/m\textsuperscript{3}; the point is the example of the text, and the band the 20 to 50~kJ of a yielding bolt (Table~\ref{tab:support-props}) on an assumed 1~m pattern. Right: displacement capacities of the example of the text; the system is limited to the 20~mm of the sprayed concrete.}
\label{fig:burst-energy}
\end{figure}

Destressing (blasting or hydraulically fracturing to move the stress
concentration away from the boundary) is used where the demand cannot be
supported. Destressing relocates the stress concentration but does not reduce
the energy stored in the surrounding rock, and its operational cost is
significant, so it is a supplement to a properly designed dynamic support
system and not a substitute for one.

\section{Monitoring in high-stress and squeezing ground}
\label{sec:high-stress-monitoring}

Monitoring in these conditions differs from routine practice in what is
measured and how quickly it must be interpreted.

In squeezing ground the quantity of interest is the convergence rate and its
trend. Convergence decelerating towards a limit indicates a system finding
equilibrium; convergence at constant or increasing rate indicates one that is
not, and it is the trend rather than the magnitude that triggers action.
Extensometers distinguish deep-seated movement of the rock mass from
loosening of the immediate skin, which determines whether longer reinforcement
would help.

In burst-prone ground microseismic monitoring is the primary tool. Arrays of
geophones locate and quantify seismic events, revealing where stress is
concentrating and where energy is being released, often well before any
visible damage. The data support decisions about sequencing, destressing and
re-entry times after blasting, and re-entry protocols based on the decay of
seismic activity are standard practice in deep mining.

Both applications require the interpretation to be available
quickly enough to act on. The observational method requires
the response time of the instruments and of their interpretation to be short
in comparison with the rate at which the ground can change
(Chapter~\ref{ch:monitoring}). A monitoring system that reports weekly in
ground that can deteriorate within a shift does not meet that requirement.

\chapterreferences
